\section{Global vacuum fidelity of local scalar read programs}
\label{sec:m1-vacuum-fidelity}

The finite original-vacuum comparison in
Section~\ref{sec:source-scalar-quantum} distinguishes a stationary reference
from a squeezed split-step baseline. On a growing population this distinction
has a global consequence: stable evolution and agreement on smooth probes do
not suffice to control the total energy produced above the original vacuum.
This section proves sharp full-spectrum thresholds and an explicit repair.
Every comparison retains the original Hamiltonian and its original vacuum.

\subsection{Exact mode error and a constructive higher-order step}

For a real symmetric positive matrix $B$ on $N$ sites, use the finite
Schr\"odinger representation with
\[
 H=\tfrac12(P^TP+Q^TBQ),\qquad [Q_i,P_j]=i\hbar\delta_{ij}.
\]
Let $\Omega$ be its normalized ground state and let
$\omega_k=\sqrt{\lambda_k(B)}$. A potential kick of signed duration $s$
sends $P$ to $P-sBQ$; an onsite kinetic drift sends $Q$ to $Q+sP$.
The kick--drift--kick step of duration $\tau$ has, in the mode coordinates
$X=\sqrt\omega Q$, $Y=P/\sqrt\omega$, matrix
\[
 S_2(z)=\begin{pmatrix}1-z^2/2&z\\-z+z^3/4&1-z^2/2\end{pmatrix},
 \qquad z=\tau\omega.
\]
Put $\rho=\sqrt[3]{2}$, $w=(2-\rho)^{-1}$ and $v=-\rho(2-\rho)^{-1}$. Use the
standard symmetric fourth-order composition
\[
 S_4(z)=S_2(wz)S_2(vz)S_2(wz).
\]
This composition is due to Yoshida, \emph{Physics Letters A} 150 (1990),
262--268, DOI \href{https://doi.org/10.1016/0375-9601(90)90092-3}
{10.1016/0375-9601(90)90092-3}. Its middle inverse step is retained;
inverse Hamiltonian gates are operations, not negative transport time.

Both matrices have the form $S=\left(\begin{smallmatrix}a&b\\c&a\end{smallmatrix}\right)$
with $a^2-bc=1$. Exact multiplication gives
\begin{align*}
 a_4&=1-z^2/2+z^4/24+(1/48+5\rho/288+\rho^2/72)z^6,\\
 b_4&=z-z^3/6-(1/36+\rho/36+\rho^2/48)z^5,\\
 c_4&=-z+z^3/6+(\rho+\rho^2)z^5/144
             -(25/1728+5\rho/432+\rho^2/108)z^7.
\end{align*}
For $0<z\le1/10$, these formulas imply $b>0$, $c<0$, $0<a<1$ and
$0.9z^2\le-bc\le1.1z^2$. Indeed $1259/1000<\rho<126/100$ bounds every
coefficient rationally, and the remaining powers obey $z^2\le1/100$.
Moreover $b_2+c_2=z^3/4$ and
\[
 b_4+c_4=d_4z^5-Dz^7,\quad
 d_4=-1/36-\rho/48-\rho^2/72\ne0,\quad
 D=25/1728+5\rho/432+\rho^2/108>0,
\]
with $z^5/16\le|b_4+c_4|\le z^5/8$ on that domain.

\begin{proposition}[Exact original-vacuum excitation]
\label{prop:m1-vacuum-excitation}
Set $\cos\theta=a$, $\sin\theta=\sqrt{-bc}>0$ and
$\chi=\sqrt{-c/b}$. After $j$ steps the original particle number in one
mode and the full-population quantities are
\begin{align}
 n_j(z)&=\frac{(b+c)^2}{-4bc}\sin^2(j\theta),
 &N_j&=\sum_k n_j(\tau\omega_k),\label{eq:m1-vacuum-number}\\
 \Delta E_j&=\hbar\sum_k\omega_k n_j(\tau\omega_k),
 &F_j&=|\langle\Omega,U^j\Omega\rangle|^2
       =\prod_k(1+n_j(\tau\omega_k))^{-1/2}.
\end{align}
In particular the second-order step produces exactly $z^6/64$ particles
in one mode after one step.
\end{proposition}
\begin{proof}
The exact power is
\[
 S^j=\begin{pmatrix}
 \cos(j\theta)&\sin(j\theta)/\chi\\
 -\chi\sin(j\theta)&\cos(j\theta)
 \end{pmatrix}.
\]
The initial covariance is $\hbar I/2$. Applying this matrix gives
$n_j=(\chi-\chi^{-1})^2\sin^2(j\theta)/4$, which is
\eqref{eq:m1-vacuum-number}. At $j=1$, $n_1=(b+c)^2/4$.
The overlap of the normalized squeezed Gaussian with the original vacuum
is $(1+n_j)^{-1/2}$ in squared modulus. Orthogonal diagonalization gives
$N$ independent real modes, so their energies add and their overlaps
multiply. Fourier multiplicities are counted in full. The quadratic
identities hold on Schwartz space; the metaplectic unitaries and bounded
effects act on the full finite-population Hilbert space.
\end{proof}

The invariant squeezed vacuum is a different reference. It carries
original-mode excess energy
$\hbar\omega(\chi+\chi^{-1}-2)/4$. Substituting it does not remove an
error relative to $H$ and $\Omega$.

\subsection{Sharp full-population thresholds}

\begin{theorem}[Global vacuum and energy fidelity]
\label{thm:m1-global-vacuum-threshold}
Let $B$ range over positive population matrices satisfying
\[
 \sigma^2=N^{-1}\operatorname{Tr}B\longrightarrow\infty,
 \qquad \lambda_{\max}(B)\le2\sigma^2.
\]
For $p=2$ or $4$, take the respective
step above, $\tau\to0$ and $\tau\sqrt{\lambda_{\max}}\le1/10$.
For any fixed $T>0$, average over every $j=1,\ldots,\lfloor T/\tau\rfloor$.
Then
\begin{align*}
 \overline{N_j}&=\Theta(N\tau^{2p}\sigma^{2p}),&
 \overline{-\log F_j}&=\Theta(N\tau^{2p}\sigma^{2p}),\\
 \overline{\Delta E_j}&=\Theta(\hbar N\tau^{2p}\sigma^{2p+1}).
\end{align*}
The corresponding upper bounds hold uniformly for every integer $j$.
Consequently uniform vacuum fidelity on a fixed nonzero time interval is
equivalent to $N\tau^{2p}\sigma^{2p}\to0$, and uniform total energy
fidelity is equivalent to $N\tau^{2p}\sigma^{2p+1}\to0$.
\end{theorem}
\begin{proof}
At least $N/3$ eigenvalues exceed $\sigma^2/2$; otherwise the mean is
less than $\sigma^2$ even if all the remaining eigenvalues are $2\sigma^2$.
Thus for every fixed $s>0$,
\[
 \frac N3(\sigma/\sqrt2)^s\le\sum_k\omega_k^s
                         \le N(\sqrt2\sigma)^s.
\]
The explicit nonzero defects above give
$c_pz^{2p}\le(b+c)^2/(-4bc)\le C_pz^{2p}$. The finite geometric sum gives
\[
 \left|\frac1J\sum_{j=1}^J\sin^2(j\theta)-\frac12\right|
 \le\frac{1}{2J\sin\theta}.
\]
For the bulk modes the right side is $O((T\sigma)^{-1})$. Using these
modes for the lower bound and all modes for the upper bound proves the
number and energy statements. The inequalities
$x/(1+x)\le\log(1+x)\le x$ and the uniform bound on each $n_j$ give
the fidelity statement. Uniform convergence implies convergence of the
time average, proving necessity. Sufficiency uses the all-time upper bound.
\end{proof}

If either scale diverges, the corresponding error is large on a positive
fraction of the times, since its mean and maximum have matching orders.
The theorem permits isolated revivals; it does not select favorable times.
When $\tau\sigma\to0$, the sharper energy average is
\[
 \overline{\Delta E_j}
 =\frac{\hbar d_p^2}{8}\tau^{2p}\sum_k\omega_k^{2p+1}(1+o(1)),
 \qquad d_2=1/4.
\]
The phase-average remainder has relative size $O((T\sigma)^{-1})$ and
the coefficient remainder has size $O((\tau\sigma)^2)$. In particular
state-norm convergence does not imply convergence of the unbounded energy:
at a constant nonzero energy scale the number scale tends to zero, leaving
rarer excitations of increasing frequency.

For the vacuum, every bounded effect has probability error at most
$\sqrt{1-F_j}$. More generally, take coherent inputs of original excess
energy at most $E_0$, supported on $\omega\le\Omega_0$ with a fixed gap
$\omega\ge\mu>0$. In dimensionless mode coordinates their displacement
norm is at most $\sqrt{2E_0/(\hbar\mu)}$. On the stated stability domain,
$|\chi-1|,|\chi^{-1}-1|\le z^p$ and $|\theta-z|\le z^{p+1}$.
For the last bound, $|a-\cos z|\le z^4/24$ for $S_2$ and
$|a-\cos z|<0.102z^6$ for $S_4$, whereas both angle sines exceed $0.94z$.
Thus comparison to the exact oscillator rotation gives, when
$\tau\Omega_0\le1/10$ and for every bounded effect,
\[
 |\Delta\Pr|\le\sqrt{1-F_j}
  +\tau^p\Omega_0^p(1+T\Omega_0)\sqrt{E_0/(\hbar\mu)}.
\]
This follows from Weyl covariance, the coherent overlap and the triangle
inequality for trace distance. It is a global observable comparison for
these preparations, not an unrestricted input-channel norm claim.

\subsection{The sparse ultraviolet scale and its repair}

On the unit periodic cube, take $N=q^3$ sites, $c=\mu=\hbar=1$ and a
symmetric nonzero integer stencil $V$ without wrapped duplicate edges. Define
\[
 d=|V|,\quad M_2=\sum_{v\in V}|v|^2,\quad
 A f(x)=\frac{6q^2}{M_2}\sum_{v\in V}(f(x)-f(x+v)),\quad B=I+A.
\]
Writing $W=6q^2d/M_2$, one has $\sigma^2=1+W$ and
$\lambda_{\max}(B)\le1+2W\le2\sigma^2$. The theorem therefore applies
to these actual graphs without a spectral success filter.

For integers $t\ge1$, set
\[
 q=2^{16t},\quad m=2^{8t},\quad K=2^{6t},\quad r=2^{7t},\quad R=mK.
\]
Let $U$ be the union of $m$ times the closed integer $K$-ball, the closed
integer $r$-ball and the signed axial lengths $2^j$, $0\le j<8t$, with
zero removed. The two balls meet only at zero. Ball Riemann sums give
\[
 d\sim\frac{4\pi}{3}r^3=\Theta(q^{21/16}),\qquad
 M_2\sim\frac{4\pi}{5}m^2K^5=\Theta(q^{23/8}),\qquad
 W\sim10q^{7/16}.
\]
The short ball dominates degree; the spaced ball dominates second moment.
Thus $\sigma=\Theta(q^{7/32})$, whereas the longest physical read is
$a=R/q=q^{-1/8}$. Its inverse is not the bulk frequency scale.

\begin{corollary}[An explicit globally faithful sparse program]
\label{cor:m1-vacuum-repair}
Use $S_4$ on $U$ with $\tau=q^{-5/8}/64$. At every level it is in the
proved stability domain, and uniformly in the step index,
\[
 \Delta E_j=O(q^{-1/32}),\qquad 1-F_j=O(q^{-1/4}).
\]
The bounded-effect comparison above follows for the stated coherent inputs.
The same tick with $S_2$ has time-averaged excess energy
$\Theta(q^{51/32})$ on each fixed nonzero time interval.
\end{corollary}
\begin{proof}
Insert $N=q^3$ and $\sigma=\Theta(q^{7/32})$ in the theorem. For all-level
stability, integer cube bounds give $d\le60r^3$ and
$M_2\ge m^2K^5/256$. The latter follows by retaining the box
$K/4\le v_1\le K/2$, $|v_2|,|v_3|\le K/4$ in the integer $K$-ball.
Thus $W\le92160q^{7/16}$ and
$\tau\sqrt{1+2W}\le\sqrt{184321}\,q^{-13/32}/64<1/10$ for $q\ge2^{16}$.
\end{proof}

\subsection{Read cost, clock scope and an exact-preservation obstruction}

Merging adjacent potential kicks, $J$ second-order steps use $J+1$ kicks
and $J$ drifts; $J$ fourth-order steps use $3J+1$ kicks and $3J$ drifts.
Every reference kick reads each site's $d$ current neighbors. The ordered
nonlocal counts are exactly $(J+1)Nd$ and $(3J+1)Nd$. The quantum potential
unitary factors into $Nd/2$ commuting two-mode phases and $N$ onsite mass
phases per kick; a drift is $N$ onsite kinetic unitaries. The classical
writer ledger records reference calculations, not cloned quantum fields.

At fixed horizon and small energy tolerance $\epsilon$, the matching bounds
give the within-method cost
\[
 \mathrm{reads}=\Theta\!\left(
 T d N^{1+1/(2p)}\sigma^{1+1/(2p)}\epsilon^{-1/(2p)}\right).
\]
The resulting powers of $q$ are $683/128$ and $1263/256$ for sparse $S_2$
and $S_4$. For the dense radius-$q^{-1/2}$ ball they are $47/8$ and $87/16$,
respectively, at the same $q$ and energy tolerance. The explicit repair uses
$\Theta(q^{79/16})$ incidences and its stated vanishing energy tolerance.
Equal $q$ is not equal total continuum accuracy. These are declared arithmetic
and gate counts, not native bit, energy or execution-time lower bounds.

Each kick expands canonical support by at most one stencil edge; each drift
is onsite. The support bound after $J$ macros is $(J+1)a$ or $(3J+1)a$.
Here $a/\tau$ diverges. The simulation duration is therefore not identified
with a physical transport clock, and inverse gates do not make it negative.

\begin{proposition}[Attained signal front and simulation-clock obstruction]
\label{prop:m1-split-clock-obstruction}
Before periodic wrapping, the canonical support radii $(j+1)a$ for $S_2$
and $(3j+1)a$ for $S_4$ are attained by nonzero signals. A faithful
implementation with complete influence speed at most $c$ must assign at
least these distances divided by $c$ as elapsed time. The simulation tick of the
explicit quantum repair cannot be this native clock at any fixed finite $c$.
\end{proposition}
\begin{proof}
In each stated stencil, $Re_1$ is the unique displacement with first
coordinate $R$. Its coefficient in $B$ is $-\alpha$, where $\alpha>0$.
For $S_4$, require $(3j+1)R<q/2$. The coefficient from initial $Q_y$ to
final $P_x$, $x-y=(3j+1)Re_1$, is exactly
\[
 -D(2A)^{j-1}\tau^{6j+1}(-\alpha)^{3j+1}\ne0,
\]
where $A=1/48+5\rho/288+\rho^2/72>0$ and $D>0$ is given above.
Cayley--Hamilton gives the lower-left entry of $S^j$ as
$U_{j-1}(a)$ times that of $S$, and $U_{j-1}$ has leading coefficient
$2^{j-1}$. The highest terms in unscaled $(Q,P)$ coordinates are
$A\tau^6B^3$ on the diagonal and $-D\tau^7B^4$ below it.
Only the path using $Re_1$ at every factor of $B^{3j+1}$ reaches that
extreme site. Lower powers cannot reach it, and the non-wrapping condition
excludes alternative periodic images. For $S_2$ the same argument gives
$(-1)^{j-1}\tau^{2j+1}(-\alpha)^{j+1}/4\ne0$ at $(j+1)Re_1$.

A local $Q$ displacement changes the final $P_x$ mean by this nonzero
coefficient. A bounded sine effect on a finite Gaussian baseline has a
nonzero response for suitable finite displacement and probe strength.
These are operational influences in the declared oscillator family, so a
complete speed bound requires the stated elapsed times. On the repair,
choose $j=\lfloor1/(8a)\rfloor$. Its support distance tends to $3/8$
without wrapping, while $j\tau=O(q^{-1/2})\to0$. No fixed finite $c$
can supply this timing identification.
\end{proof}

The LC reduction bounds every admitted native influence, including small
ones. The proposition excludes this clock attachment; it does not assert
that the quantum approximation fails. Other operation grammars and
implementations are outside this obstruction. Local ultraviolet interventions are
not the fixed-band preparations of the global comparison. Six full 32-site
ray controls verify the extreme coefficients and zeros outside support for
both methods at $j=1,2,3$.

There is also an exact obstruction. For a scalar polynomial symplectic
update $S(B)$ preserving the original vacuum, its mode matrix preserves
$\operatorname{diag}(\lambda,1)$. Comparing the energy and symplectic
inverses gives $d=a$, $c=-\lambda b$ and
\[
 a(\lambda)^2+\lambda b(\lambda)^2=1.
\]
If there are more distinct eigenvalues than this polynomial's degree, it is
an identity. The two nonconstant terms have opposite degree parity, so their
leading coefficients cannot cancel: $b=0$ and $a=\pm1$. A consistent
near-identity step is the identity. The stated families have unboundedly
many distinct eigenvalues: along $(k,0,0)$, for
$0\le k\le\lfloor q/(4R)\rfloor$, each nonzero axial projection contributes
a strictly increasing $1-\cos(2\pi k v_1/q)$, with angle at most $\pi/2$.
Thus no fixed number of scalar kick/drift stages nontrivially preserves the
original vacuum exactly on all refinements. This does not exclude
nonpolynomial flows, growing-degree interpolation or other operation grammars.

\paragraph{Evidence and scope.}
The independent packet replays four complete periodic spectra, 512
population observations, 24 full spatial covariance comparisons, 702 forward
and 702 inverse scalar writes, and 2,268 nonlocal reference reads. Exact
cubic arithmetic checks the signed composition and its stability domain.
The $q=65536$ moment certificate uses integer columns; it does not execute
that full population. Ten Lean reductions have transitive standard-axiom
audits. The universal spectral and asymptotic arguments are analytic.
This supplies a global quantum approximation and its explicit cost, not an
all-level A1--A3 source, finite twelve-port quantum implementation, interacting
matter law or physical clock. Existing physical premises are unchanged.
